% Options for packages loaded elsewhere
\PassOptionsToPackage{unicode}{hyperref}
\PassOptionsToPackage{hyphens}{url}
\documentclass[
]{article}
\usepackage{xcolor}
\usepackage{amsmath,amssymb}
\setcounter{secnumdepth}{-\maxdimen} % remove section numbering
\usepackage{iftex}
\ifPDFTeX
  \usepackage[T1]{fontenc}
  \usepackage[utf8]{inputenc}
  \usepackage{textcomp} % provide euro and other symbols
\else % if luatex or xetex
  \usepackage{unicode-math} % this also loads fontspec
  \defaultfontfeatures{Scale=MatchLowercase}
  \defaultfontfeatures[\rmfamily]{Ligatures=TeX,Scale=1}
\fi
\usepackage{lmodern}
\ifPDFTeX\else
  % xetex/luatex font selection
\fi
% Use upquote if available, for straight quotes in verbatim environments
\IfFileExists{upquote.sty}{\usepackage{upquote}}{}
\IfFileExists{microtype.sty}{% use microtype if available
  \usepackage[]{microtype}
  \UseMicrotypeSet[protrusion]{basicmath} % disable protrusion for tt fonts
}{}
\makeatletter
\@ifundefined{KOMAClassName}{% if non-KOMA class
  \IfFileExists{parskip.sty}{%
    \usepackage{parskip}
  }{% else
    \setlength{\parindent}{0pt}
    \setlength{\parskip}{6pt plus 2pt minus 1pt}}
}{% if KOMA class
  \KOMAoptions{parskip=half}}
\makeatother
\usepackage{longtable,booktabs,array}
\usepackage{caption}
\captionsetup[table]{skip=6pt}
\newcounter{none} % for unnumbered tables
\usepackage{calc} % for calculating minipage widths
% Correct order of tables after \paragraph or \subparagraph
\usepackage{etoolbox}
\makeatletter
\patchcmd\longtable{\par}{\if@noskipsec\mbox{}\fi\par}{}{}
\makeatother
% Allow footnotes in longtable head/foot
\IfFileExists{footnotehyper.sty}{\usepackage{footnotehyper}}{\usepackage{footnote}}
\makesavenoteenv{longtable}
\usepackage{graphicx}
\makeatletter
\newsavebox\pandoc@box
\newcommand*\pandocbounded[1]{% scales image to fit in text height/width
  \sbox\pandoc@box{#1}%
  \Gscale@div\@tempa{\textheight}{\dimexpr\ht\pandoc@box+\dp\pandoc@box\relax}%
  \Gscale@div\@tempb{\linewidth}{\wd\pandoc@box}%
  \ifdim\@tempb\p@<\@tempa\p@\let\@tempa\@tempb\fi% select the smaller of both
  \ifdim\@tempa\p@<\p@\scalebox{\@tempa}{\usebox\pandoc@box}%
  \else\usebox{\pandoc@box}%
  \fi%
}
% Set default figure placement to htbp
\def\fps@figure{htbp}
\makeatother
\setlength{\emergencystretch}{3em} % prevent overfull lines
\providecommand{\tightlist}{%
  \setlength{\itemsep}{0pt}\setlength{\parskip}{0pt}}
\usepackage{bookmark}
\IfFileExists{xurl.sty}{\usepackage{xurl}}{} % add URL line breaks if available
\urlstyle{same}
% fallback for those not using the hyperref driver hyperxmp:
\makeatletter
\@ifundefined{xmpquote}{\newcommand{\xmpquote}[1]{#1}}{}
\makeatother
\hypersetup{
  hidelinks,
  pdfcreator={LaTeX via pandoc}}

\author{}
\date{}

\begin{document}

\section{Bab 7: Kegagalan Tata Kelola: D3TLH Dalam Sistem
Perizinan}\label{bab-7-kegagalan-tata-kelola-d3tlh-dalam-sistem-perizinan}

\textbf{CELIOS --- Center of Economic and Law Studies}

\emph{Evaluasi terhadap kelumpuhan supremasi hukum di mana instrumen
perlindungan ekologis direduksi menjadi stempel administratif.}

\begin{center}\rule{0.5\linewidth}{0.5pt}\end{center}

\subsection{Metodologi Pendekatan \& Pertanyaan
Kritis}\label{metodologi-pendekatan-pertanyaan-kritis}

\textbf{Membaca Hubungan Antara Hasil D3TLH dan Keputusan Perizinan
Aktual}

Analisis di halaman ini menggunakan \textbf{Matriks Analisis Crosstab}
untuk menjawab 3 pertanyaan fundamental: 1. \textbf{Apakah D3TLH
digunakan sebagai dasar keputusan?} 2. \textbf{Apakah D3TLH bersifat
mengikat atau hanya rekomendasi?} 3. \textbf{Apakah D3TLH dapat
diabaikan secara prosedural?}

\emph{Kerangka Pengujian:} Menyilangkan Data Fase Status Ekologis (Aman
/ Tertekan / Kritis) dengan Data Empiris Keputusan Izin yang benar-benar
diterbitkan negara.

\begin{center}\rule{0.5\linewidth}{0.5pt}\end{center}

Secara hukum, dokumen tata ruang dan Daya Dukung Daya Tampung Lingkungan
Hidup (D3TLH) dirancang sebagai \textbf{veto absolut} negara---sebuah
instrumen yang memiliki otoritas untuk menolak izin investasi baru
manakala beban lingkungan hidup di suatu ekoregion telah melampaui
kapasitas pemulihannya. Namun, matriks forensik yang membedah
\emph{timeline} penerbitan izin di semenanjung Sulawesi membongkar
skandal \textbf{Regulatory Capture}. Negara, melalui aparatus
perizinannya, terbukti mendelegitimasi instrumen penyelamatannya
sendiri. Dokumen AMDAL dan D3TLH diposisikan bukan sebagai garis merah
hukum (\emph{redline}), melainkan sekadar ornamen pelengkap persyaratan
untuk merestui perluasan tambang dan pembangunan masif PLTU captive.

\begin{center}\rule{0.5\linewidth}{0.5pt}\end{center}

\subsection{Temuan Kunci Tata Kelola}\label{temuan-kunci-tata-kelola}

{\def\LTcaptype{none} % do not increment counter
\begin{longtable}[]{@{}
  >{\raggedright\arraybackslash}p{(\linewidth - 4\tabcolsep) * \real{0.3333}}
  >{\raggedright\arraybackslash}p{(\linewidth - 4\tabcolsep) * \real{0.3333}}
  >{\raggedright\arraybackslash}p{(\linewidth - 4\tabcolsep) * \real{0.3333}}@{}}
\toprule\noalign{}
\begin{minipage}[b]{\linewidth}\raggedright
Indikator
\end{minipage} & \begin{minipage}[b]{\linewidth}\raggedright
Status / Nilai
\end{minipage} & \begin{minipage}[b]{\linewidth}\raggedright
Keterangan
\end{minipage} \\
\midrule\noalign{}
\endhead
\bottomrule\noalign{}
\endlastfoot
\textbf{Skor Tata Kelola (Veto)} & \textbf{9.8 / 10} & Tingkat keparahan
kegagalan fungsi D3TLH sebagai pembatas. \textbf{Vonis: Regulatory
Capture.} (Indeks Komposit CELIOS) \\
\textbf{Otoritas D3TLH} & \textbf{Nihil} & Instrumen lingkungan terbukti
tidak mengikat dan gagal membatasi laju obral perizinan secara nyata di
lapangan. \\
\textbf{Status Penegakan Hukum} & \textbf{Pembiaran} & Masifnya temuan
perusahaan beroperasi secara ilegal tanpa dikenakan sanksi pencabutan
izin operasi (Impunitas massal). \\
\end{longtable}
}

\begin{center}\rule{0.5\linewidth}{0.5pt}\end{center}

\subsection{7.1 Pembuktian Empiris: Status Ekologis vs Penerbitan
Izin}\label{pembuktian-empiris-status-ekologis-vs-penerbitan-izin}

\textbf{Metode: Spatial Overlay \& Crosstabulation (ESDM x GFW)}

\subsubsection{Metodologi: Evaluasi Kepatuhan D3TLH Berdasarkan Data
Historis}\label{metodologi-evaluasi-kepatuhan-d3tlh-berdasarkan-data-historis}

\textbf{Metode Analisis:} Sub-bab ini menggunakan agregasi berbasis
aturan (\emph{Rule-based Categorization}) untuk membedah ketidaksesuaian
antara status kerusakan lingkungan dengan keputusan administratif
perizinan.

\begin{enumerate}
\def\labelenumi{\arabic{enumi}.}
\tightlist
\item
  \textbf{Model Evaluasi Pelanggaran (\emph{Compliance Modeling}):}

  \begin{itemize}
  \tightlist
  \item
    \textbf{Kategorisasi Status (Binning):} Nilai kerusakan lingkungan
    absolut dibagi ke dalam tiga kelas menggunakan distribusi
    \emph{percentile}: Aman (≤33\%), Tertekan (33-66\%), dan Kritis
    (\textgreater66\%).
  \item
    \textbf{Kuantifikasi Pelanggaran D3TLH:} Mengidentifikasi secara
    kuantitatif apakah pemerintah tetap mengobral Izin Usaha
    Pertambangan (IUP) baru pada wilayah-wilayah yang secara empiris
    terbukti telah berada di fase `Kritis'.
  \end{itemize}
\item
  \textbf{Kalkulasi/Formula Pengolahan:}

  \begin{itemize}
  \tightlist
  \item
    \texttt{Ambang\_Kritis\ =\ Percentile(Deforestasi,\ 0.66)}
  \item
    \texttt{Total\_Izin\_Ilegal\_Ekologis\ =\ SUM(IUP\_Baru)\ WHERE\ Status\_D3TLH\ =\ \textquotesingle{}Kritis\textquotesingle{}}
  \end{itemize}
\item
  \textbf{Variabel \& Fitur Data:}

  \begin{itemize}
  \tightlist
  \item
    \textbf{Variabel Konteks Lingkungan:}
    \texttt{Total\_Deforestasi\_Ha} atau
    \texttt{Deforestasi\_Driver\_Komoditas...} (sebagai basis status
    wilayah)
  \item
    \textbf{Variabel Keputusan Aktor:} \texttt{Jumlah\_Izin\_Baru} dan
    \texttt{Total\_Luas\_Konsesi\_Baru\_Ha}
  \end{itemize}
\item
  \textbf{Dataset \& File:}

  \begin{itemize}
  \tightlist
  \item
    \texttt{data/processed/sulawesi\_izin\_baru\_per\_tahun.csv} dan
    \texttt{sulawesi\_gfw\_master\_1\_dekade\_2014\_2023.csv}
  \end{itemize}
\end{enumerate}

\begin{center}\rule{0.5\linewidth}{0.5pt}\end{center}

Dalih ``sudah sesuai aturan'' sering kali digunakan oleh pemerintah
daerah maupun pusat untuk membenarkan penerbitan Izin Usaha Pertambangan
(IUP) yang masif di Pulau Sulawesi. Namun, jika Daya Tampung dan Daya
Dukung Lingkungan Hidup (D3TLH) benar-benar dijadikan sebagai instrumen
pencegahan dan batas pengaman ekologis (\emph{ecological safeguard}),
maka secara logika hukum dan sains, laju penerbitan izin baru wajib
dihentikan seketika saat sebuah wilayah telah menunjukkan gejala
kerusakan ekstrem, seperti deforestasi besar-besaran. Kenyataannya, data
spasial menunjukkan sebuah anomali fatal: grafik penerbitan izin justru
melonjak eksponensial tepat di tahun-tahun ketika hilangnya tutupan
pohon sekunder dan primer mencapai titik kritis.

Kami menyilangkan data deforestasi tahunan dari \emph{Global Forest
Watch} (GFW) dengan riwayat penerbitan izin pertambangan dari
\emph{Minerba One Data Indonesia} (MODI) Kementerian ESDM. Hasil irisan
tersebut memperlihatkan pola \emph{business-as-usual} yang sangat
agresif. Sebagai contoh, ketika jutaan hektar hutan di Sulawesi Tengah
dan Tenggara ditebang secara brutal pada periode 2014 hingga 2023 untuk
membuka jalan bagi kawasan industri (smelter) dan infrastruktur
pendukungnya, pemerintah secara bersamaan terus mencetak ratusan Izin
Usaha Pertambangan baru di provinsi-provinsi yang sama, tanpa
mempertimbangkan bahwa kapasitas hidrologis dan daya dukung lanskap di
wilayah tersebut sudah ambruk total.

Fakta bahwa perizinan konsesi tetap mengalir deras di zona-zona merah
ini membuktikan kegagalan tata kelola yang sistemik. D3TLH, Analisis
Mengenai Dampak Lingkungan (AMDAL), maupun instrumen Kajian Lingkungan
Hidup Strategis (KLHS) telah direduksi menjadi sekadar stempel
administratif formalitas belaka. Alih-alih menjadi rem darurat bagi
kerusakan, dokumen-dokumen lingkungan tersebut ditundukkan oleh tekanan
investasi dan target pertumbuhan produksi nikel nasional. Kegagalan
otoritas dalam menggunakan data riil kerusakan lapangan sebagai dasar
penghentian operasi atau moratorium izin ini pada akhirnya mempercepat
laju krisis ekologis yang kini harus ditanggung sepenuhnya oleh
masyarakat pesisir, petani, dan nelayan yang kehilangan ruang hidup
mereka. Kehancuran ini bukan lagi kecelakaan, melainkan hasil dari
kebijakan pembiaran (\emph{state-sponsored ecocide}).

Matriks Statistik di bawah ini digunakan untuk membuktikan secara
kuantitatif jejak kelumpuhan tata kelola (\emph{Governance Failure})
tersebut.

\subsubsection{Konklusi Analisis Kepatuhan D3TLH Berdasarkan Data
Historis}\label{konklusi-analisis-kepatuhan-d3tlh-berdasarkan-data-historis}

Tabel pembuktian di bawah ini mengukur akumulasi penerbitan izin pada
rentang waktu ketika wilayah berstatus Aman, Tertekan, hingga Kritis
secara aktual. Apabila pada status Kritis izin masih diterbitkan, hal
tersebut secara matematis mendiskualifikasi D3TLH sebagai instrumen
perlindungan lingkungan.

{\def\LTcaptype{none} % do not increment counter
\begin{longtable}[]{@{}
  >{\raggedright\arraybackslash}p{(\linewidth - 8\tabcolsep) * \real{0.2000}}
  >{\raggedright\arraybackslash}p{(\linewidth - 8\tabcolsep) * \real{0.2000}}
  >{\raggedright\arraybackslash}p{(\linewidth - 8\tabcolsep) * \real{0.2000}}
  >{\raggedright\arraybackslash}p{(\linewidth - 8\tabcolsep) * \real{0.2000}}
  >{\raggedright\arraybackslash}p{(\linewidth - 8\tabcolsep) * \real{0.2000}}@{}}
\toprule\noalign{}
\begin{minipage}[b]{\linewidth}\raggedright
Status Daya Dukung
\end{minipage} & \begin{minipage}[b]{\linewidth}\raggedright
Kondisi Kerusakan Hutan
\end{minipage} & \begin{minipage}[b]{\linewidth}\raggedright
Seharusnya (Menurut Aturan)
\end{minipage} & \begin{minipage}[b]{\linewidth}\raggedright
Kenyataan di Lapangan
\end{minipage} & \begin{minipage}[b]{\linewidth}\raggedright
Kesimpulan Tata Kelola
\end{minipage} \\
\midrule\noalign{}
\endhead
\bottomrule\noalign{}
\endlastfoot
\textbf{Aman} & Ringan (Hilang 2,649 - 14,542 Ha) & Wajar diterbitkan
izin & 39 Izin Baru Keluar & Normal (Sesuai Aturan) \\
\textbf{Tertekan} & Sedang (Hilang 14,568 - 36,808 Ha) & Izin mulai
direm/dibatasi & 81 Izin Baru Keluar & Anomali (Lampu Kuning) \\
\textbf{Kritis} & Sangat Parah (Hilang 39,633 - 154,169 Ha) &
\textbf{Moratorium / Izin Dilarang!} & \textbf{260 Izin Baru Keluar}
(Termasuk luasan 447,683 Ha) & \textbf{BUKTI PELANGGARAN FATAL} \\
\end{longtable}
}

\begin{quote}
\textbf{Temuan Target:} - \textbf{Daya dukung tidak menjadi pembatas
nyata} (Terdapat \textbf{260 Izin Baru} yang tetap sengaja diterbitkan
ketika lingkungan sudah terkonfirmasi berada di fase Kritis). -
\textbf{Keputusan perizinan tetap dominan secara politis} meski D3TLH
menunjukkan kondisi kritis (Bukti \emph{Regulatory Capture} mutlak).
\end{quote}

\begin{center}\rule{0.5\linewidth}{0.5pt}\end{center}

\subsection{7.2 Tabrakan Hukum: Impunitas dan Pembiaran Operasi
Ilegal}\label{tabrakan-hukum-impunitas-dan-pembiaran-operasi-ilegal}

\textbf{Impunitas Korporasi dan Pembiaran Konflik Struktural di Sektor
Ekstraktif}

\textbf{Metode: Thematic Coding \& Analisis Kasus (LSM / KPA / Tanah
Kita)}

\subsubsection{Metodologi: Pemetaan Impunitas
Korporasi}\label{metodologi-pemetaan-impunitas-korporasi}

\textbf{Metode Analisis:} Sub-bab ini menggunakan agregasi pelaporan
berbasis insiden (\emph{Incident-based Reporting Aggregation}) untuk
mengukur tingkat pembiaran penegakan hukum (impunitas).

\begin{enumerate}
\def\labelenumi{\arabic{enumi}.}
\tightlist
\item
  \textbf{Model Penelusuran Anomali Hukum:}

  \begin{itemize}
  \tightlist
  \item
    \textbf{Kasus Rekam Jejak:} Menyaring dan mengklasifikasikan
    database konflik sengketa lahan, pelanggaran HAM, dan kasus operasi
    ilegal tanpa izin di level tapak.
  \item
    \textbf{Pemetaan Pembiaran (\emph{State Omission}):} Menghitung
    total volume agregat di mana korporasi yang terbukti bermasalah
    secara hukum tetap dipertahankan keberadaan operasinya oleh aparatur
    negara.
  \end{itemize}
\item
  \textbf{Kalkulasi/Formula Pengolahan:}

  \begin{itemize}
  \tightlist
  \item
    \texttt{Total\_Kasus\_Impunitas\ =\ COUNT(Judul\_Kasus)}
  \item
    \texttt{Volume\_Pembiaran\_Sektoral\ =\ SUM(Kasus)\ GROUP\ BY\ Sektor}
  \end{itemize}
\item
  \textbf{Variabel \& Fitur Data:}

  \begin{itemize}
  \tightlist
  \item
    \textbf{Atribut Laporan:} \texttt{Provinsi}, \texttt{Sektor},
    \texttt{Judul\_Kasus}, \texttt{Deskripsi\_Singkat}
  \end{itemize}
\item
  \textbf{Dataset \& File:}

  \begin{itemize}
  \tightlist
  \item
    \texttt{data/processed/sulawesi\_konflik\_hukum.csv}
  \end{itemize}
\end{enumerate}

\begin{center}\rule{0.5\linewidth}{0.5pt}\end{center}

Konsep Daya Tampung dan Daya Dukung Lingkungan Hidup (D3TLH) secara
esensial tidak hanya mengukur batas fisik daya tahan ekosistem darat,
air, dan udara, tetapi juga mengukur sejauh mana ruang hidup sosial
masyarakat mampu menoleransi masuknya investasi skala raksasa. Namun,
realitas di lapangan menunjukkan bahwa ambang batas sosial ini telah
dilanggar secara brutal dan sistematis. Berdasarkan kompilasi data
laporan masyarakat sipil, Jaringan Advokasi, serta Konsorsium Pembaruan
Agraria (KPA), kami menemukan bahwa ekspansi industri ekstraktif di
Sulawesi---mulai dari pertambangan nikel, perkebunan monokultur, hingga
pembangunan kawasan industri strategis nasional (PSN)---selalu
beriringan dengan letusan konflik agraria yang memakan korban di pihak
masyarakat akar rumput.

Pembiaran atas pelanggaran hukum (impunitas) terlihat sangat telanjang.
Perusahaan-perusahaan yang terbukti merebut lahan tanpa persetujuan
bebas, diinformasikan di awal, dan tanpa paksaan atau \emph{Free, Prior
and Informed Consent} (FPIC) dari masyarakat adat maupun penduduk lokal,
tetap dibiarkan beroperasi dengan leluasa. Dalam berbagai kasus,
korporasi tambang maupun perkebunan yang status izin (HGU/IUP)-nya
tumpang tindih dengan tanah ulayat, atau bahkan beroperasi di luar
konsesi yang sah, tidak pernah mendapatkan sanksi administratif berupa
pencabutan izin. Aparat penegak hukum justru kerap digunakan sebagai
instrumen represi untuk membungkam protest warga, mengkriminalisasi
petani yang mempertahankan kebunnya, serta menyingkirkan masyarakat yang
bermukim di sekitar area lingkar tambang.

Absennya sanksi tegas dari Kementerian Lingkungan Hidup, Kementerian
ESDM, maupun aparat hukum mencerminkan adanya perlindungan politik yang
kuat terhadap operasi bisnis kotor ini. Konflik tenurial yang terjadi di
Morowali, Konawe, hingga Wawonii membuktikan bahwa instrumen hukum yang
seharusnya melindungi rakyat telah dibajak demi memfasilitasi akumulasi
modal. Pembiaran ini menciptakan tabrakan hukum di mana regulasi
perlindungan lingkungan hidup dan hak asasi manusia kalah secara
sistemik oleh undang-undang sektoral pro-investasi seperti UU Cipta
Kerja dan UU Minerba revisi terbaru. Kegagalan menegakkan hukum ini pada
akhirnya mengonfirmasi bahwa negara secara sadar menoleransi kekerasan
struktural dan kerusakan lingkungan sebagai harga yang wajar demi
pencapaian target industri hilirisasi, mengubah wilayah-wilayah krisis
tersebut menjadi zona pengorbanan (\emph{sacrifice zones}).

{\def\LTcaptype{none} % do not increment counter
\begin{longtable}[]{@{}
  >{\raggedright\arraybackslash}p{(\linewidth - 4\tabcolsep) * \real{0.3333}}
  >{\raggedright\arraybackslash}p{(\linewidth - 4\tabcolsep) * \real{0.3333}}
  >{\raggedright\arraybackslash}p{(\linewidth - 4\tabcolsep) * \real{0.3333}}@{}}
\toprule\noalign{}
\begin{minipage}[b]{\linewidth}\raggedright
Indikator Impunitas
\end{minipage} & \begin{minipage}[b]{\linewidth}\raggedright
Nilai
\end{minipage} & \begin{minipage}[b]{\linewidth}\raggedright
Keterangan
\end{minipage} \\
\midrule\noalign{}
\endhead
\bottomrule\noalign{}
\endlastfoot
\textbf{Total Kasus Konflik/Pelanggaran Dibiarkan} & \textbf{32 Kasus} &
Bukti Impunitas Hukum di Sulawesi (KPA/LSM) \\
\end{longtable}
}

\begin{figure}
\centering
\pandocbounded{\includegraphics[keepaspectratio,alt={Timeline Letusan Konflik Agraria}]{visuals_bab7/chart_7_2_konflik_hukum.png}}
\caption{Timeline Letusan Konflik Agraria}
\end{figure}

\begin{center}\rule{0.5\linewidth}{0.5pt}\end{center}

\subsection{7.3 Inkonsistensi Iklim: Karpet Merah PLTU
Captive}\label{inkonsistensi-iklim-karpet-merah-pltu-captive}

\textbf{Paradoks Hilirisasi Hijau dan Karpet Merah untuk PLTU Batubara
Captive}

\textbf{Metode: Penyaringan Agregat Dataset Eksternal (Global Coal Plant
Tracker GEM)}

\subsubsection{Metodologi: Agregasi Beban Karbon PLTU
Captive}\label{metodologi-agregasi-beban-karbon-pltu-captive}

\textbf{Metode Analisis:} Sub-bab ini menggunakan inventarisasi agregat
kuantitatif (\emph{Quantitative Inventory Aggregation}) dari database
global PLTU batubara independen (\emph{captive}).

\begin{enumerate}
\def\labelenumi{\arabic{enumi}.}
\tightlist
\item
  \textbf{Model Ekstraksi Kapasitas Fosil:}

  \begin{itemize}
  \tightlist
  \item
    \textbf{Isolasi Regional:} Melakukan pemfilteran data inventaris
    energi kotor (PLTU) yang berlokasi secara presisi di kawasan
    industri strategis pulau Sulawesi.
  \item
    \textbf{Kuantifikasi Kontradiksi Karbon:} Menghitung total jumlah
    \emph{unit} pembangkit dan agregat luaran listrik kotor (dalam
    satuan Megawatt) yang dibangun secara masif demi menopang pabrik
    pemurnian nikel, yang notabene dipromosikan sebagai proyek energi
    ramah lingkungan.
  \end{itemize}
\item
  \textbf{Kalkulasi/Formula Pengolahan:}

  \begin{itemize}
  \tightlist
  \item
    \texttt{Total\_Beban\_Karbon\ =\ SUM(Capacity\_MW)\ GROUP\ BY\ Provinsi}
  \item
    \texttt{Total\_Infrastruktur\_Kotor\ =\ COUNT(Unit\_PLTU)}
  \end{itemize}
\item
  \textbf{Variabel \& Fitur Data:}

  \begin{itemize}
  \tightlist
  \item
    \textbf{Spesifikasi Pembangkit:} \texttt{Capacity\ (MW)},
    \texttt{Start\ year}, \texttt{Provinsi\ (Subnational\ unit)}
  \end{itemize}
\item
  \textbf{Dataset \& File:}

  \begin{itemize}
  \tightlist
  \item
    \texttt{data/processed/sulawesi\_pltu\_captive.csv}
  \end{itemize}
\end{enumerate}

\begin{center}\rule{0.5\linewidth}{0.5pt}\end{center}

Di panggung negosiasi iklim internasional, pemerintah secara retoris
selalu mempromosikan komitmen Indonesia dalam transisi energi menuju
target \emph{Net Zero Emission} serta mengkampanyekan keberhasilan
program ``Hilirisasi Hijau'' untuk mendukung rantai pasok kendaraan
listrik (EV) global. Namun, ironi terbesar dan paling memalukan dari
narasi transisi energi tersebut terbentang sangat nyata di pulau
Sulawesi. Demi memenuhi kebutuhan listrik raksasa untuk operasional
pabrik pemurnian (\emph{smelter}) nikel, pemerintah memberikan karpet
merah perizinan untuk pembangunan puluhan Pembangkit Listrik Tenaga Uap
(PLTU) Batubara \emph{Captive}. Pembangkit-pembangkit kotor ini dibangun
secara mandiri di dalam kawasan industri dan berdiri di luar sistem
jaringan (grid) PLN nasional, sehingga keberadaannya sering kali luput
dari skema pembatasan emisi atau target penghentian batubara (\emph{coal
phase-out}).

Ekstraksi data terbaru dari \emph{Global Coal Plant Tracker} (GEM)
membongkar skalanya yang masif dan mengerikan. Kami menemukan bahwa
terdapat puluhan unit PLTU \emph{Captive}---baik yang sudah beroperasi
secara aktif, sedang dalam tahap konstruksi, maupun yang telah
mendapatkan izin (\emph{permitted/announced})---yang tersebar di
sentra-sentra industri di Sulawesi Tengah, Sulawesi Tenggara, dan
provinsi lainnya. Total kapasitas kotor pembangkitan dari puluhan PLTU
batubara ini mencapai belasan ribu Megawatt. Ini berarti, proses
pengolahan material nikel yang diklaim sebagai bahan baku energi bersih
masa depan (baterai kendaraan listrik) justru diproduksi dengan membakar
jutaan metrik ton batu bara setiap tahunnya, menyemburkan polutan
mematikan seperti SO2, NOx, dan PM2.5 langsung ke ruang udara pemukiman
warga sekitar.

Inkonsistensi kebijakan ini jelas menunjukkan kegagalan tata kelola
iklim yang paling telanjang. Transisi energi dalam konteks hilirisasi
nikel di Indonesia tidak lebih dari sekadar \emph{greenwashing}, di mana
perlindungan daya dukung udara dikorbankan demi profitabilitas ekonomi.
Emisi Gas Rumah Kaca (GRK) dari Sulawesi justru meledak secara dramatis
berkat kehadiran PLTU \emph{Captive} ini. Dengan terus memberikan izin
baru bagi pembangkit listrik berbahan bakar fosil terburuk di bumi ini,
pemerintah tidak hanya mengkhianati komitmen penurunan emisinya sendiri,
tetapi juga menjerumuskan masyarakat lokal ke dalam kondisi kerentanan
kesehatan pernapasan akut yang tak terpulihkan.

{\def\LTcaptype{none} % do not increment counter
\begin{longtable}[]{@{}
  >{\raggedright\arraybackslash}p{(\linewidth - 4\tabcolsep) * \real{0.3333}}
  >{\raggedright\arraybackslash}p{(\linewidth - 4\tabcolsep) * \real{0.3333}}
  >{\raggedright\arraybackslash}p{(\linewidth - 4\tabcolsep) * \real{0.3333}}@{}}
\toprule\noalign{}
\begin{minipage}[b]{\linewidth}\raggedright
Indikator PLTU Captive
\end{minipage} & \begin{minipage}[b]{\linewidth}\raggedright
Nilai
\end{minipage} & \begin{minipage}[b]{\linewidth}\raggedright
Keterangan
\end{minipage} \\
\midrule\noalign{}
\endhead
\bottomrule\noalign{}
\endlastfoot
\textbf{Total Unit PLTU Captive} & \textbf{67 Unit} &
Beroperasi/Dibangun/Direncanakan Khusus Kawasan Sulawesi \\
\textbf{Total Kapasitas Pembangkitan Kotor} & \textbf{12,245 MW} &
Pembangkit Fosil Penopang Smelter \\
\end{longtable}
}

\begin{figure}
\centering
\pandocbounded{\includegraphics[keepaspectratio,alt={Timeline Pertumbuhan Kapasitas PLTU Captive}]{visuals_bab7/chart_7_3_pltu_captive_timeline.png}}
\caption{Timeline Pertumbuhan Kapasitas PLTU Captive}
\end{figure}

\end{document}
